\documentclass[11pt,a4paper]{article}
\usepackage[margin=25mm]{geometry}
\usepackage{amsmath,amssymb,amsthm,booktabs}
\newtheorem{lemma}{Lemma}\newtheorem{theorem}[lemma]{Theorem}
\title{Four O'Nan map counts from one small group}
\author{Human Compression project}\date{}
\begin{document}\maketitle
We determine the ordered map counts of types $(7,16)$, $(14,16)$
and $(16,28)$, together with $(8,28)$, for $G=O'N$. We reuse the complete subgroup roster,
the exact character evaluations and one explicit group of order 1344.
The argument counts overlaps of maximal subgroups without multiplying
paths in the subgroup-class diagram.

\section{The possible generated subgroups}
The two cyclic order-16 classes are rows 510 and 512. In their
55-row upward closure, exactly the following rows have order divisible by 7:
\[
1,2,3,4,5,7,9,10,21,53,54,65,66,93,95,115,116,117,118.
\]
Every subgroup containing elements of orders 16 and $q\in\{7,14,28\}$
lies in this list. Row 21 is itself maximal and has no additional
maximal overgroup. Its qualifying pairs are retained in the full subtraction.
The reader verifies the cyclic classes, the complete closure, every order
and the two recorded cover relations directly in the fixed roster.

Rows 2 and 3 are the two maximal $L_3(7):2$ classes; rows 7 and 21
are maximal of orders 161280 and 10752, respectively. The direct incidence data give unique actual immediate
overgroups along the chains
\[
93\longrightarrow53\longrightarrow9\longrightarrow4\longrightarrow2,
\qquad95\longrightarrow54\longrightarrow10\longrightarrow5\longrightarrow3.
\]
For each arrow $K\to L$, the class length of $L$ times the number of
$K$ subgroups in one $L$ equals the class length of $K$. Double counting
incidences gives one actual overgroup. Each subgroup in these chains
therefore has a unique maximal overgroup.

Rows 65 and 66 have order 1344 and only immediate-overgroup classes
$\{2,7\}$ and $\{3,7\}$, respectively. Their direct data are
\begin{center}\begin{tabular}{rrrr}\toprule
smaller row&larger row&smaller per larger&larger class length\\\midrule
65&2&2793&122760\\65&7&120&2857239\\
66&3&2793&122760\\66&7&120&2857239\\\bottomrule
\end{tabular}\end{center}
Both smaller class lengths are $342868680$, equal to every product
in the last two columns. Thus each has exactly two maximal overgroups.
Rows 115--118 have order 672 and immediate-overgroup classes
$\{9,65\}$, $\{10,66\}$, $\{65\}$ and $\{66\}$, respectively.
The direct incidences give each exactly one overgroup of order 1344.
These four rows require the centralizer argument below.

\section{A literal model of the group of order 1344}
Use the already acquired model of $L_3(7):2$ on the 57 projective
points and 57 projective lines of $\mathbb F_7^3$. A determinant-one
matrix acts by $v\mapsto vg$ on points and $v\mapsto vg^{-T}$ on
lines. The graph involution interchanges the two halves with the same
coordinates. This realization and the identification of both O'Nan
maximal classes with this abstract group are the same named classical
premises used in the earlier $(3,12)$ proof.

For $A\in\operatorname{GL}_2(7)$, set
\[
 \iota(A)=\operatorname{diag}((\det A)^{-1},A)\in\operatorname{SL}_3(7).
\]
The projective kernel consists of $A=\lambda I_2$ with $\lambda^3=1$,
so its image has order $2016/3=672$. The graph involution normalizes
this image, since $\iota(A)^{-T}=\iota(A^{-T})$, and lies outside it
by its action on the two halves. Thus the extension $H$ has order 1344.
The finite certificate independently verifies the exact set of 1344
permutations and its generation by the images of
\[
 \begin{pmatrix}1&1\\0&1\end{pmatrix},\qquad
 \begin{pmatrix}0&-1\\1&0\end{pmatrix},\qquad
 \begin{pmatrix}3&0\\0&1\end{pmatrix}
\]
and the graph involution. It checks their projective membership and
finds 336 elements of order 16, 96 of order 28, and exactly
\[
 x_H(7,16)=x_H(14,16)=2688,\qquad x_H(16,28)=5376
\]
ordered pairs $(r,s)$ of the indicated orders with $|rs|=2$. The complete element list and the $(16,28)$ pair
indices are retained. A separate reader reconstructs the generated
group, computes all element orders and tests every eligible pair.
It also verifies that the eighth powers of all 336 order-16 elements
are the same nonidentity involution $z$, and that $z$ commutes with
the four generators. Hence $z\in Z(H)$.

This model is identified by containment and order, not by a structure
label alone. Map it into a maximal subgroup in row 2. Its elements of
orders 16 and 7 place its image in the candidate list above. Its order
1344 leaves only rows 65 and 66; row 66 cannot be contained in row 2,
since its only proper overgroups are the distinct maximal classes
3 and 7. The image therefore has row 65. Applying the same argument
in row 3 gives row 66. 
\section{The remaining overgroups}
\begin{lemma}\label{centralizers}
Let $M=\operatorname{PSL}_3(7):\langle\tau\rangle$, where
$\tau$ is inverse transpose. An inner involution has centralizer of
order 1344 in $M$. The centralizer of an outer involution has no
element of order 16.
\end{lemma}
\begin{proof}
An inner involution lifts, after multiplication by a scalar cube
root of unity, to an involution of $\operatorname{SL}_3(7)$. It is
conjugate to $\operatorname{diag}(-1,-1,1)$. A projective centralizer
cannot send this matrix to a nontrivial scalar multiple: their traces
would differ. Its centralizer in $\operatorname{SL}_3(7)$ is the block
group $\operatorname{GL}_2(7)$, of order 2016. Quotienting by the
scalar centre of order three gives 672. Inverse transpose centralizes
the diagonal involution, so the centralizer in $M$ has order 1344.

Write an outer involution as $z=[A]\tau$ with $A\in
\operatorname{SL}_3(7)$. The involution condition gives
$AA^{-T}=\zeta I$ with $\zeta^3=1$, or $A=\zeta A^T$.
Transposing again gives $\zeta^2=1$, hence $\zeta=1$ and $A$ is
symmetric. The inner centralizer is the projective similitude group
of this nondegenerate three-dimensional symmetric form. Every
similitude multiplier $\mu$ is a square, since its determinant
satisfies $(\det h)^2=\mu^3$. Thus each projective similitude has a
unique representative in the special orthogonal group: first rescale
to an isometry, then choose the sign making its determinant one.
These representatives belong to $\operatorname{SL}_3(7)$, so no
additional restriction arises from passing to $\operatorname{PSL}_3(7)$.

A semisimple element of 2-power order in this odd-dimensional special
orthogonal group has eigenvalues $1,\lambda,\lambda^{-1}$. Frobenius
either fixes the two nontrivial eigenvalues or exchanges them. Thus
their orders divide 6 or 8, respectively, and the element's 2-power
order divides 8. Finally $C_M(z)=C_{\operatorname{PSL}_3(7)}(z)
\times\langle z\rangle$, so adjoining the outer involution still
introduces no element of order 16.
\end{proof}

Let $K$ be in one of rows 115--118, and let $H_0$ be its unique
order-1344 overgroup. Since $K$ contains an order-16 element, the
literal eighth-power property puts the corresponding central
involution $z$ of $H_0$ in $K$. It is central in $K$.
If a maximal subgroup $M\cong L_3(7):2$ contains $K$, then
$K\leq C_M(z)$. Lemma~\ref{centralizers} excludes the outer case
because $K$ contains an element of order 16. Thus $C_M(z)$ has order
1344 and contains $K$ with index two. It is an immediate overgroup
of $K$, and the direct incidence calculation forces it to equal $H_0$.
Therefore every such $M$ contains the same $H_0$.

A maximal overgroup in row 7 also contains $H_0$. An immediate
overgroup of $K$ inside it cannot be of row 9 or 10: their order 32928
is divisible by $7^2$, whereas 161280 is not. The remaining possibility
is $H_0$. The complete upward lists show these are all maximal-overgroup
classes. Thus each $K$ has exactly the same two maximal overgroups as $H_0$.

\section{Exact subtraction}
Write $x_L=x_L(p,q)$ for the number of ordered pairs of the specified
orders with product of order two. A nongenerating pair generates a
subgroup in the candidate list. Outside rows 65, 66 and 115--118, that
subgroup has exactly one maximal overgroup. In these six rows it has
exactly two, and lies in a unique subgroup in rows 65 or 66. Conversely,
every qualifying pair in an order-1344 subgroup generates one of the
six classes, by the candidate list and subgroup orders. The overlap
pairs are therefore partitioned by the actual order-1344 subgroups.
Consequently
\[
 e_G(p,q)=x_G(p,q)-\sum_M[G:M]x_M(p,q)
             +2(342868680)x_H(p,q).
\]
The sum includes all thirteen maximal conjugacy classes. Dividing by
$|\operatorname{Aut}G|=921631011840$ gives
\begin{center}\begin{tabular}{rrrr}\toprule
type&earlier lower bound&correction&exact count\\\midrule
$(7,16)$&7446&2&7448\\
$(14,16)$&12502&2&12504\\
$(16,28)$&25312&4&25316\\\bottomrule
\end{tabular}\end{center}
There is no additional duality division. All three values agree with
Connor--Leemans, Table~1. \section{The same model for type $(8,28)$}
Intersecting the upward closures of cyclic order-eight and order-28
subgroups gives exactly
\[
                    1,2,3,7,8,65,66,67,112,113.               \tag{*}
\]
The direct incidence calculation gives one actual parent on each
of the additional edges
\[
8\to7,\quad67\to7,\quad112\to8,\quad112\to66,
\quad113\to8,\quad113\to65.
\]
The complete immediate-parent lists for rows 112 and 113 are exactly
$\{8,66\}$ and $\{8,65\}$; rows 8 and 67 have only parent class 7.
Thus rows 8 and 67 have one maximal overgroup.

Fix $K$ in row 113 and its unique parent $H_0$ in row 65. Let $M$
be the unique maximal overgroup of $H_0$ in row 7. It has a unique
normal subgroup $E$ of index two, in row 8. Since $H_0$ is maximal
in $M$ and distinct from $E$, the intersection $H_0\cap E$ has
order 672. The complete list of subgroups of that order in $H_0$
is one each in rows 113, 115 and 117. The latter two contain
an element of order 16, whereas row 8 does not, as the complete
cyclic order-16 closure shows. Consequently $H_0\cap E=K$.
This identifies $E$ as the unique row-8 parent of $K$, and therefore
joins the two routes to the same maximal subgroup in row 7.
Any projective maximal overgroup of $K$ must contain $H_0$:
the other possible immediate parent, of order 80640, cannot lie
in a group of order 3753792. The same argument, with rows 112,
66, 116 and 118, proves the other case. Each of rows 112 and 113
therefore has exactly the same two maximal overgroups as its
unique parent of order 1344.

It follows from (*) that a qualifying pair is counted twice among
maximal overgroups precisely when its generated subgroup lies in
rows 65, 66, 112 or 113. Such a pair lies in one actual subgroup of
row 65 or 66. Conversely every qualifying pair in either of these
groups generates one of those four rows: the other candidates in
(*) have incompatible orders or are different classes of the same
order. Thus the same overlap formula applies.

The unchanged 1344-element set contains 168 elements of order eight
and 96 of order 28. Testing all their ordered pairs gives
$x_H(8,28)=1344$. The independent reader repeats this complete test;
there is no second model acquisition. Hence the correction is
\[
 \frac{2(342868680)(1344)}{921631011840}=1,
 \qquad \#\text{map classes of type }(8,28)=6275+1=6276.
\]
This is also the fixed Table~1 value.

The classification and character inputs are
borrowed; the contribution is the small literal model and the overlap
proof. No chirality, principal-refutation or conjecture credit is claimed.
\begin{thebibliography}{9}
\bibitem{CL} T. Connor and D. Leemans, \emph{Algorithmic enumeration
of regular maps}, Ars Math. Contemp. 10 (2016), 211--222.
\bibitem{roster} T. Connor and D. Leemans, \emph{The subgroup lattice
of $O'N$}, fixed 23-page edition; native source joins accompany the
literal subgroup certificate.
\end{thebibliography}
\end{document}
